\begin{tikzpicture}[
    >=Stealth,
    font=\sffamily,
    node distance=1.8cm and 2.5cm,
    % Estilo de las cajas de entidades
    entity/.style={
        rectangle,
        rounded corners=4pt,
        draw=cyan!70!black,
        fill=cyan!20!black,
        text=white,
        thick,
        align=center,
        inner sep=10pt,
        minimum width=3.8cm,
        minimum height=1.4cm
    },
    % Estilo del régimen de restricciones
    regime/.style={
        rectangle,
        rounded corners=6pt,
        draw=orange!80!black,
        fill=orange!20!black,
        text=white,
        thick,
        align=center,
        inner sep=12pt,
        minimum width=9.5cm
    },
    % Estilos de líneas
    downward/.style={->, very thick, color=orange!80!black},
    upward/.style={->, thick, dashed, color=teal!70!black},
    mutual/.style={<->, very thick, color=cyan!70!black}
]

    % Régimen de restricciones (nivel superior)
    \node[regime] (CR) {
        \textbf{Régimen de Restricciones (\textit{Constraint-Regimes})} \\
        \small (Scripts culturales, lenguaje, instituciones, infraestructura material)
    };

    % Entidades / Agentes (nivel inferior)
    \node[entity, below left=2.2cm and 0.2cm of CR] (EA) {
        \textbf{Entidad / Agente A} \\
        \footnotesize Modelo generativo (AIF)
    };
    
    \node[entity, below right=2.2cm and 0.2cm of CR] (EB) {
        \textbf{Entidad / Agente B} \\
        \footnotesize Modelo generativo (AIF)
    };

    % 1. Interacción horizontal (multidireccional entre agentes)
    \draw[mutual] (EA) -- (EB) 
        node[midway, below=3pt, align=center, font=\scriptsize, text=white] {
            Intercambio bidireccional /\\
            acoplamiento mutuo
        };

    % 2. Influencia vertical descendente (Top-down: constriñe y modula)
    \draw[downward] ([xshift=-1.8cm]CR.south) -- (EA.north)
        node[midway, left=2pt, align=right, font=\scriptsize, text=orange!80!black] {
            Ejerce límites \\ y modula el régimen (\textit{Poder})
        };
        
    \draw[downward] ([xshift=1.8cm]CR.south) -- (EB.north)
        node[midway, right=2pt, align=left, font=\scriptsize, text=orange!80!black] {
            Refuerzan o \\ alteran los scripts (\textit{Poder})
        };

    % 3. Influencia vertical ascendente / retroactiva (Bottom-up: sedimentación o cambio)
    \draw[upward] (EA.north west) to[out=110, in=205] ([xshift=-3.5cm]CR.south)
        ;
    \node[anchor=south east, xshift=-0.25cm, yshift=0.55cm,
        align=right, font=\scriptsize, text=teal!70!black] at (EA.north west) {
        (top-down)
    };

    \draw[upward] (EB.north east) to[out=70, in=335] ([xshift=3.5cm]CR.south)
        ;
    \node[anchor=south west, xshift=0.25cm, yshift=0.55cm,
        align=left, font=\scriptsize, text=teal!70!black] at (EB.north east) {
        (Retroalimentación \\ bottom-up)
    };

    % Fondo que enmarca la totalidad como un sistema interactivo
    \begin{scope}[on background layer]
        \node[
            draw=gray!30,
            dashed,
            rounded corners=8pt,
            fill=black,
            fit=(CR) (EA) (EB),
            inner sep=14pt,
            label={[anchor=south east, font=\itshape\tiny, text=gray!70]south east:Sistema dinámico no lineal}
        ] {};
    \end{scope}

\end{tikzpicture}